%% ===========================================================================
\section{Conclusion}
\label{sec:conclusion}
%% ===========================================================================

We studied a hybrid architecture in which inexpensive extractors propose candidate values and a judge selects one of them or abstains, and we asked whether a small decision model can play the judge better than large open-weight language models. On three benchmarks, two languages and single- and multi-turn inputs, a fine-tuned Laya judge with fewer than half a billion parameters was more accurate than LLM judges of 27 to 180 billion parameters, 3 to 6 times faster per request and 5 to 38 times faster in batch throughput, on a fraction of one GPU. The main reason is abstention: the LLM judges choose well when the right value is present, but too often answer when the user expressed nothing. After calibration, the small judge's confidences also rank its errors far better, which lets a system decide most requests on its own at a few percent risk and ask back on the rest; the LLM judges cannot do so on the hardest dataset. Two design lessons emerge. Robustness to the shape of the option lists comes from training-time augmentation only when the augmentation reproduces the confusions of real proposers. And cross-lingual transfer remains imperfect: zero-shot transfer from English costs 0.04--0.09 F1, and joint English--French training, which helps with oracle lists, becomes unstable with real candidates unless the checkpoint is selected on matching data. With real candidates, the bottleneck moves to the proposers: a fine-tuned sequence tagger alone outperforms our pipeline when values are literal spans of the utterance, while the pipeline wins when values come from the dialogue history or must be mapped to canonical forms. The judge is therefore best seen not as a competitor of extraction models but as the decision layer on top of them.

Beyond slot filling, these results support a general design principle for agentic and dialogue systems: when a decision is frequent, well defined and can be supervised, a compact model trained with a proper scoring rule on the exact decision format, and calibrated on held-out data, can be more accurate, cheaper and more trustworthy than prompting a general model. The next steps are to use the tagger as one of the proposers, to learn the decision policy from asymmetric costs and from user feedback with reinforcement learning, at the level of single requests and of whole dialogues, and to study few-shot LLM judges, the two-step decision, unseen domains and robustness to schema changes.
